\documentclass[12pt]{article}
\usepackage[utf8]{inputenc}
\usepackage[T1]{fontenc}
\usepackage{hyperref}
\title{Tagish Language Resources for K-12 Curriculum Development}
\author{Compiled by Claude Sonnet 5 (claude-sonnet-5) for Andr\'as Kornai}
\date{}
\begin{document}
\maketitle

\section{Introduction}
Tagish (ISO 639-3 tgx) is an Athabaskan language of the southern Yukon, Canada. It has no remaining first-language speakers. Revitalization is led by the Carcross/Tagish First Nation, with the Yukon Native Language Centre (YNLC) and others. The sources below are mostly from 2004 to 2018, so current programme status needs confirming with the community.

\section{Primary Educational Resources}
\begin{itemize}
\item A Carleton University news story reports that the YNLC, a department of the Council of Yukon First Nations, is digitizing its teaching resources for Yukon languages including Tagish. The material is nearly 10,000 pages of literacy sessions, language lessons and storybooks, free to search and download on the YNLC site: \url{https://carleton.ca/news/story/recovering-yukon-languages/}
\item The Council of Ministers of Education, Canada, has a document on the Yukon FirstVoices Aboriginal Languages Project: \url{https://cmec.ca/docs/Aboriginal/YK-FirstVoicesAboriginalLanguagesProject.en.pdf}
\end{itemize}

\section{Linguistic and Reference Materials}
\begin{itemize}
\item Endangered Languages Project resource list for Tagish: \url{https://endangeredlanguages.com/node/37821/resources}
\item An academic paper on the Tagish FirstVoices Project, ``New Technologies and Contested Ideologies'': \url{https://paperguide.ai/papers/0c1deff6-76db-4b30-9ef0-d8a3d49632d5/}
\end{itemize}

\section{Religious and Community Resources}
A Council of Yukon First Nations release on Yukon language work: \url{https://cyfn.ca/?p=3202}

\section{Digital and Audio Resources}
\begin{itemize}
\item The Carcross/Tagish First Nation built a Tagish FirstVoices website with sound files, photographs, videos and text. It was funded under a 2004 agreement, as described in a Yukon government release: \url{https://bac-lac.wayback.archive-it.org/web/20080206205124/http://www.gov.yk.ca/news/2005/files/05-288.pdf}
\item National Research Council project in which community language workers across Yukon's 14 First Nations created video recordings of Elders speaking their languages: \url{https://nrc.canada.ca/en/research-development/research-collaboration/programs/project-create-video-recordings-yukon-indigenous-languages}
\end{itemize}

\section{Curriculum Development Context}
Tagish is a revitalization case. The YNLC lesson and storybook collection, the FirstVoices site and the Elder video recordings are the best starting points. Whether Tagish-specific materials sit inside the YNLC collection should be checked with the YNLC and the Carcross/Tagish First Nation before any curriculum work.
\end{document}